% Generated draft: requires mathematical and source review.
\documentclass[11pt]{article}
\usepackage[margin=25mm]{geometry}
\usepackage{fontspec}
\setmainfont{Latin Modern Roman}
\usepackage{amsmath,amssymb,mathtools,unicode-math}
\setmathfont{Latin Modern Math}
\usepackage{enumitem,needspace,xurl,hyperref}
\setlength{\parindent}{0pt}
\setlength{\parskip}{5pt}
\setlength{\emergencystretch}{3em}
\newcommand{\N}{\mathbb{N}}
\newcommand{\Z}{\mathbb{Z}}
\newcommand{\Q}{\mathbb{Q}}
\newcommand{\R}{\mathbb{R}}
\newcommand{\C}{\mathbb{C}}
\newcommand{\F}{\mathbb{F}}
\newcommand{\A}{\mathbb{A}}
\newcommand{\PP}{\mathbb{P}}
\newcommand{\E}{\mathbb{E}}
\newcommand{\eps}{\varepsilon}
\newcommand{\dd}{\,\mathrm{d}}
\begin{document}
\section*{UnsolvedMath ID: 20000338}
\textbf{Source problem code:} catalogue entry
\subsection{Definitions and mathematical statement}
\textbf{Original request.} Find natural loci in moduli spaces of surfaces whose mutual intersections would be interesting to study. No particular moduli space, surface class, field, notion of naturalness, or desired intersection invariant is specified.

\textbf{Attached unreviewed worked proposal.} The aid considers the plane-octic branch chart for Gorenstein stable surfaces with $K_X^2=2$ and $\chi(\mathcal O_X)=4$. In the usual double-cover interpretation the branch is an octic in $\PP^2$; the coefficient field and characteristic hypotheses are not stated in the summary. Write $V_j=H^0(\PP^2,\mathcal O(j))$ and use the ambient parameter space $\PP(V_8)$. The line-factor locus is the image of $\PP(V_1)\times\PP(V_7)$ under multiplication, and the conic-factor locus is the image of $\PP(V_2)\times\PP(V_6)$. The aid gives their codimensions as $7$ and $12$.

Near a generic branch $lqf=0$, where $l,q,f$ have degrees $1,2,5$, respectively, the aid reports that these loci meet along a smooth codimension-$17$ locus. The intersection is clean: the tangent space to their intersection is the intersection of their tangent spaces. Its excess rank is $7+12-17=2$. The generic open conditions include a transverse line-conic intersection $Z=L\cap Q$ of length two and no vanishing of $f$ at those points; complete stability and transversality conditions need source review.

For the explicit bundle notation, let
\[T=\PP(V_1)\times\PP(V_2)\times\PP(V_5),\]
let $a,b,c$ be the respective hyperplane classes, and let $p:\mathcal Z\to T$ be the universal length-two intersection of the line and conic on the open set where it is finite flat and transverse. The notation $\mathcal O_{\mathcal Z}(8)$ means the restriction of $\mathcal O_{\PP^2}(8)$ from the plane factor. The aid's excess-bundle formula is
\[\mathcal E=\mathcal O_T(1,1,1)\otimes p_*\mathcal O_{\mathcal Z}(8),\qquad\operatorname{rank}\mathcal E=2,\]
\[c_2(\mathcal E)=136a^2+108ab+36b^2+17ac+9bc+c^2.\]
It identifies the excess normal fiber, with the corresponding tautological line factor understood on the projective parameter space, with $H^0(\mathcal O_{L\cap Q}(8))$ and canonically with the common smoothing space of the two $A_1$ surface singularities lying above $L\cap Q$. These are the two singularities associated with that intersection, not a claim that no other branch intersections produce singularities.

The reported degree of the projective closure of the $(1,2,5)$-factorized locus in the octic parameter space is $18{,}648{,}630$. The aid explicitly limits these statements to a local/open branch-chart result: neither this degree nor the displayed second Chern class is asserted to be a global KSBA moduli intersection number. Descent through projective coordinate changes and extension across the boundary are not supplied.
\subsection{Short English statement}
The source asks for interesting pairs of natural surface-moduli loci. The aid proposes the line-factor and conic-factor loci for octic double-cover branches, with a clean excess-two intersection, a computed rank-two excess bundle, and a projective parameter-space degree. These claims remain local to the branch chart and are not global moduli intersection numbers.
\subsection{Sources}
\begin{itemize}
\item Local UnsolvedMath registry, permanent ID 20000338, source code catalogue entry; \texttt{aim-algebraic-geometry-notes.json}, notes[337], tag \texttt{problem}.
\item AIM workshop \emph{Compact moduli spaces and birational geometry}, Problem 4.4: \url{https://aimath.org/WWN/birational/birational.pdf}. No individual proposer, remarks or literature references are supplied. No PDF was consulted.
\item Attached machine-generated research attempt 1, result \texttt{partial\_result}, status \texttt{partially\_solved}, rigor label \texttt{proved}, novelty confidence low. It is not peer reviewed and requires independent verification before citation. No external citation for its excess-intersection calculations is supplied.
\end{itemize}
\end{document}
